\section{Markov Mechanism Check}\label{app:v9-markov}

The synthetic Markov benchmark is the controlled mechanism check for
the local-supervision argument. The oracle is known in closed form,
the sufficient state is small, and node-level labels can be drawn at
designed propensities. Two claims are tested here. First, a flat
Fourier neural operator (FNO) trained on root labels recovers the
oracle on whole documents, so the operator family is rich enough to
represent the underlying tree aggregation. Second, when the same
problem is spread over a tree of leaves, root accuracy is steady
across a broad band of leaf sizes. The Manifesto and HLL appendices
make the parallel statements for richer oracles.

\subsection{Data-Generating Process}\label{app:v9-markov-dgp}

Each synthetic document is a fixed-length token sequence generated by
a hidden Markov chain over $K$ registers. The chain picks a register;
the register emits an observed token from its own palette of
synonyms. Register $A$ emits \emph{grief}, \emph{mourn}, \emph{tears};
register $B$ emits \emph{plot}, \emph{spy}, \emph{ploy}; and so on.
The learner sees tokens, not hidden registers. The oracle is the
number of register transitions in the hidden sequence.

\begin{figure}[t]
\centering
\includegraphics[width=\linewidth]{markov_scene_arc_hamlet.pdf}
\caption{\textbf{Boundary correction on Hamlet's Act 3 through
Ophelia's funeral.} Twelve consecutive scenes (3.1 through 5.1) are
colored by their dominant register from the five-register set
Mourning, Intrigue, Madness, Action, Political. Three same-register
runs are visible in the strip: 3.2--3.3 (Intrigue), 4.1--4.2--4.3
(Political), and 4.4--4.5 (Mourning). Leaves $L_1,\ldots,L_4$ group
three consecutive scenes; each leaf summary $S_i = g(L_i)$ stores
$(\text{first register}, \text{last register}, \text{count})$. Each
internal node adds $c_L + c_R + \One\{L.\mathrm{last}\neq
R.\mathrm{first}\}$, and the root returns the exact transition count
($=7$). Two scene-boundary flips sit between leaves
(Intrigue$\to$Action at 3.3$\to$3.4 and Intrigue$\to$Mourning at
4.7$\to$5.1). At the root merge between $L_2$ and $L_3$ the boundary
correction correctly does \emph{not} fire ($\mathrm{Political}$
vs.\ $\mathrm{Political}$), because the run that ends $L_2$
continues into $L_3$ and the leaf summaries retain their first and
last register. The Hamlet labels are a reader-facing skin on the
synthetic registers; every number reported below is computed on the
synthetic palette.}
\label{fig:v9-markov-registers}
\end{figure}

\subsection{Sufficient State}\label{app:v9-markov-state}

For any span, the sufficient state is
\[
(\mathrm{count},\mathrm{first},\mathrm{last}).
\]
The count records internal transitions; the first and last fields
carry the boundary information used by parent merges. The merge rule
is
\[
(c_L,a_L,b_L)\otimes(c_R,a_R,b_R)
=
(c_L+c_R+\One\{b_L\ne a_R\},a_L,b_R).
\]
The state is associative and ordered. Dropping the boundary fields
breaks the merge: the parent loses the information needed to decide
whether a transition crosses the child boundary.

This is the smallest setting that exposes the local-law issue. A
parent merge needs the right boundary of the left child and the left
boundary of the right child. When either child drops that boundary
state, the root count shifts even when each child summary looks
locally adequate. Figure~\ref{fig:v9-markov-registers} shows the
boundary-correction pattern on the Hamlet skin.

\subsection{Flat-FNO Baseline}\label{app:v9-markov-fno}

The flat baseline is an FNO trained from root labels alone over the
full document. On the recoverable scope, the flat FNO recovers the
oracle within a small approximation gap that tightens with training
data. Multi-seed runs at $1{,}024$ training documents place the gap
at the readout-bottleneck level expected for the FNO width and modes
(visible as the leftmost cells in
Figure~\ref{fig:v9-markov-leaf-size}). The tree model below is asked
to match or improve on this baseline once the same root supervision
is split into leaf-merge calls.

\subsection{Leaf-Size Sweep}\label{app:v9-markov-leaf-size}

Figure~\ref{fig:v9-markov-leaf-size} reports the per-cell tree
gain over the flat-FNO baseline across leaf size and the share of
documents that receive root labels. The recoverable scope (left)
shows a broad band of leaf sizes where the tree matches or improves
on the flat baseline at full supervision and remains competitive when
only $10\%$ of documents receive root labels. The structural scope
(right) is a harder DGP with more registers and more expected shifts;
the band is narrower, and very small leaves expose the
boundary-state failure mode predicted by the merge rule in
Section~\ref{app:v9-markov-state}.

\begin{figure}[t]
\centering
\begin{subfigure}[b]{0.48\textwidth}
    \centering
    \includegraphics[width=\linewidth]{markov_leaf_size_fixed_recoverable.pdf}
    \caption{Recoverable scope.}
    \label{fig:v9-markov-leaf-size-recoverable}
\end{subfigure}
\hfill
\begin{subfigure}[b]{0.48\textwidth}
    \centering
    \includegraphics[width=\linewidth]{markov_leaf_size_fixed_structural.pdf}
    \caption{Structural scope.}
    \label{fig:v9-markov-leaf-size-structural}
\end{subfigure}
\caption{\textbf{Tree-supervision gain over the flat FNO baseline
across leaf size and supervision budget.} Recoverable (left) and
structural (right) scopes on the synthetic Markov chain. The bolded
colorbar tick marks the average flat-FNO root MAE for the panel, so
cells above the tick improve on the flat baseline. The recoverable
panel shows a wide safe band of leaf sizes; the structural panel
shows the same dominance band on a narrower interior of leaf
sizes.}
\label{fig:v9-markov-leaf-size}
\end{figure}

The takeaway from these heatmaps is the steadiness band rather than
a single best cell. A range of leaf sizes preserves root accuracy
once the local laws hold on the realized leaves and merges; the band
is wider on the recoverable scope and narrower on the structural
scope.

\subsection{Supervision Allocation}\label{app:v9-markov-allocation}

The same DGP supports a controlled budget accounting. A
full-document label on a $128$-token sequence costs $128$ token
observations. A label on a $16$-token leaf costs $16$ token
observations, so eight such leaf labels have the same budget mass as
one root label. The root-label share ranges from $100\%$ to $10\%$;
the remaining mass is reallocated to local labels under three
policies. \emph{Leaf-only} spends the local budget on leaf count
labels. \emph{Depth-equal} splits local labels across non-root
depths. \emph{Balanced-node} spreads labels across non-root nodes.
Each policy assigns node supervision at known propensities, matching
the audit design in Section~\ref{sec:v8-audit-certificate}.

Figures~\ref{fig:v9-markov-simple-leaf-mass} and
\ref{fig:v9-markov-hard-leaf-mass} show the reallocation curves at
fixed token budget. The recoverable case has a broad safe band of
leaf sizes; the structural case is harder, and the tree still
dominates the flat baseline across the shown budgets.
Figure~\ref{fig:v9-markov-budget-split} isolates the supervision-object
shift: the total token budget is held fixed, the horizontal axis
changes the fraction spent on root labels versus local labels, and
the reallocation curves stay nearly flat from full root supervision
down to roughly $40\%$ root supervision before degrading.

\begin{figure}[t]
    \centering
    \includegraphics[width=0.95\linewidth]{markov_simple_leaf_mass.png}
    \caption{\textbf{Recoverable Markov supervision allocation.}
    Each panel fixes a root-label share. The green curve uses root
    labels only. The blue, purple, and rose curves keep the $100\%$
    token budget and reallocate missing root-label mass to local
    labels under the three policies. The amber curve is the flat FNO
    baseline. Lower MAE is better.}
    \label{fig:v9-markov-simple-leaf-mass}
\end{figure}

\begin{figure}[t]
    \centering
    \includegraphics[width=0.95\linewidth]{markov_hard_leaf_mass.png}
    \caption{\textbf{Structural Markov supervision allocation.}
    The structural regime repeats the same protocol with more
    registers and more expected shifts. Local supervision remains
    effective, and very small leaves expose the boundary-state
    failure mode.}
    \label{fig:v9-markov-hard-leaf-mass}
\end{figure}

\begin{figure}[t]
    \centering
    \includegraphics[width=0.95\linewidth]{markov_budget_split.png}
    \caption{\textbf{Splitting a fixed token budget between root and
    local labels.} The green star is the all-root reference. The
    green line spends only the reduced root-label budget. The blue,
    purple, and rose curves keep the total budget fixed by assigning
    the remaining mass to local labels. The flat region shows the
    range where node supervision preserves root accuracy at the same
    total budget.}
    \label{fig:v9-markov-budget-split}
\end{figure}

\subsection{Scope}\label{app:v9-markov-scope}

The Markov experiment supports a narrow methodological claim. With a
known oracle and a known sufficient state, an FNO trained on root
labels recovers the underlying aggregation, and root accuracy is
steady across a broad band of leaf sizes once the local laws hold on
the realized leaves and merges. Local labels become substitutes for
root labels when they are randomized over the relevant tree
population and measure the same oracle-preservation property; their
value is statistical as well as operational, since they provide
training signal for the compression operators and design-based
estimates of local-law distortion.

The Manifesto/Benoit experiment in the main text observes
full-document root targets over the sampled corpus and reserves
node-level observations for the stronger quasi-sentence audit path.
The classical-sketch appendix (Section~\ref{app:v8-sketch-comparisons})
carries the parallel recovery story for a randomized state with a
published error budget.
